\documentclass[10pt,a4paper]{article}
\usepackage[margin=2.2cm]{geometry}
\usepackage{fontspec}
\setmonofont{DejaVu Sans Mono}[Scale=0.84]
\newfontfamily\unifont{DejaVu Serif}[Scale=0.9]
\usepackage{amsmath,amssymb,mathtools}
\setcounter{MaxMatrixCols}{20}
\usepackage{longtable,booktabs,array}
\usepackage{fancyvrb}
\usepackage[dvipsnames]{xcolor}
\usepackage[colorlinks=true,linkcolor=NavyBlue,urlcolor=NavyBlue]{hyperref}
\usepackage{microtype}
\usepackage{needspace}
\usepackage{tocloft}\setlength{\cftbeforesecskip}{0.45em}
\setlength{\parindent}{0pt}\setlength{\parskip}{0.55em}
\setlength{\emergencystretch}{8em}
\makeatletter
\renewcommand\section{\@startsection{section}{1}{\z@}{-3.5ex \@plus -1ex \@minus -.2ex}{2.3ex \@plus.2ex}{\normalfont\Large\bfseries\raggedright}}
\renewcommand\subsection{\@startsection{subsection}{2}{\z@}{-3.25ex\@plus -1ex \@minus -.2ex}{1.5ex \@plus .2ex}{\normalfont\large\bfseries\raggedright}}
\renewcommand\subsubsection{\@startsection{subsubsection}{3}{\z@}{-3.25ex\@plus -1ex \@minus -.2ex}{1.5ex \@plus .2ex}{\normalfont\normalsize\bfseries\raggedright}}
\makeatother
\renewcommand{\arraystretch}{1.15}
\allowdisplaybreaks
\begin{document}
\title{The radion of 7+1 spacetime: provenance\\[0.4em]\large The task, the author’s note, the decisions, the inputs, the data, the derivations and the verification record of the document on the radion}
\author{Repository: github.com/once-ere/7plus1 — prepared for Nash\\ Notebook: 7+1.nb (Wolfram Language 15.0.1)\\ Document generated by tools/make\_docs.py on 2026-10-28}
\date{}
\maketitle
\begin{abstract}
This record states what was asked, how it was interpreted, which note of the author was checked and how, every decision taken (among them the new section 26 of 7+1.nb, which computes every worked example), which results of the earlier sections and which data were used and where they come from, how each result was derived, and how it was checked: the 33 checks of section 26 (33 True), the 31 boolean checks of the independent Python program tools/ref\_radion.py (31 True), the comparison of the 15 groups of numbers printed by both programs, the check that the earlier sections of the notebook are unchanged, the program versions and the SHA-256 hashes of the files.
\end{abstract}
\tableofcontents
\Needspace{12\baselineskip}\section*{1. The task}\addcontentsline{toc}{section}{1. The task}
On 28 October 2026 the author asked for the following (summarized; the original wording is the author’s private working material and is not reproduced):

\begin{enumerate}
\item Define and explain the "auxiliary torus" and the "world brane" precisely and correctly, each in its own documentation, symbol and provenance files (Markdown, LaTeX, PDF), for students, with examples, well organized; push them to the repository.
\item Then, in a second request: read and understand a note of the author on the radion (a PDF in the workspace, see section 2), and use its ideas to verify and clarify the radion and brane documents in preparation; create examples in which the four extra space dimensions are (1) a 4-sphere and (2) a 4-torus.
\end{enumerate}
The standing rules of the project apply: every documentation file has a companion provenance file and a symbol file, in Markdown and LaTeX, compiled to PDF; no file refers the reader to another documentation, provenance or symbol file; every equation that is created or used is written in the file that uses it. The repository’s .gitignore is respected.

\textbf{Interpretation.} The radion is the field that the torus and brane documents use (the size $b$ of the auxiliary torus, its coupling to the brane); it gets its own set of three documents, in which every idea of the note is stated precisely, checked by computation, and corrected where the checks require it, and every idea is worked out for the two extra spaces $T^4$ and $S^4$. The torus and brane documents gain the corresponding comparisons. "Verify" is read as: every statement of the note that can be computed is computed twice, by the notebook and by an independent program; statements that rest on data (collider bounds) are marked as quoted. The coordinates of 7+1.nb are kept: $X=(x_0,\dots,x_7)$, $x_0$ the time (the direction of the octonion unit $e_0=1$), $x_1,x_2,x_3$ ordinary space, $x_4,\dots,x_7$ the four extra space-like coordinates.

\Needspace{8\baselineskip}\subsection*{1.1 The signature: the author’s explanation of 27 October 2026}\addcontentsline{toc}{subsection}{1.1 The signature: the author’s explanation of 27 October 2026}
\textbf{Source.} On 27 October 2026 the author of the project explained that the xAct list $\{7,1,0\}$ declares seven space-like dimensions and one time-like dimension, while $\mathrm{diag}(+1,-1,\dots,-1)$ is simply the squares of the octonion basis, and asked that this be made clear in all documentation, examples and provenance records. \textbf{What it corrected.} The version of 7+1.nb pushed in the commits 52d8729 to dc2f178 of github.com/once-ere/7plus1 declared the xAct metric of section 12 with \texttt{DefMetric[\{1, 7, 0\}, …]}, reading the list as the signs of $\eta$; since then section 12 declares \texttt{DefMetric[\{7, 1, 0\}, g71[-a, -b], CD, …]}, exactly as the author’s notebook Einstein-Lovelock-7+1-Nash.nb does. \textbf{What was added.} Checks in 7+1.nb — section 4.7: $\eta$ is the squares of the octonion basis, $e_Ae_A=\eta_{AA}e_0$ and $\mathrm{Re}(e_Ae_B)=\eta_{AB}$; section 12.2: \texttt{SignatureOfMetric[g71]} $=\{7,1,0\}$ with \texttt{SignDetOfMetric[g71]} $=-1$, $\mathrm{Re}(e_Ae_B)=\eta_{AB}$, and $\mathrm{Re}(x\,x)=\eta_{AB}x^Ax^B$; section 17: the flat metric $\mathrm{diag}(+1,-1,-1,-1)$ of the world brane is the squares of the basis $e_0,e_1,e_2,e_3$ of the quaternion subalgebra — and the same statements in the Python programs \texttt{tools/\allowbreak{}ref71.py}, \texttt{tools/\allowbreak{}ref\_\allowbreak{}gravity.py} and \texttt{tools/\allowbreak{}ref\_\allowbreak{}sm.py}; the explanation itself in the text of 7+1.nb (sections 1.2, 1.4, 4.2, 4.7, 7, 11, 12.2, 13, 15.1, 17) and in every document, symbol file and provenance file of the repository. \textbf{What did not change.} No result. Compared output by output with the backup below, every output of 7+1.nb is the same, except the declared list itself, the new checks and the check counts (and, as in every run, the name of the private context and the measured times); in particular every component of every Lovelock tensor is unchanged. The components follow from the component matrix of the metric, and the sign of the determinant, \texttt{SignDetOfMetric}, is $-1$ for both lists. \textbf{Backup.} The evaluated notebook as it was before is kept in the workspace as \texttt{backups/\allowbreak{}7+1\_\allowbreak{}2026-10-27\_\allowbreak{}before\_\allowbreak{}signature.nb} (the folder backups/ is excluded from the repository by .gitignore), 1034401 bytes, SHA-256 \texttt{06094248\allowbreak{}8caf4f45\allowbreak{}120ffd37\allowbreak{}be0123e9\allowbreak{}6992a31c\allowbreak{}d7fdbde0\allowbreak{}9ee4dc4c\allowbreak{}ed4f5d3a}; \texttt{tools/\allowbreak{}verify\_\allowbreak{}7p1\_\allowbreak{}nb.py} compares every input, text and section cell of 7+1.nb with it and finds every cell present, each one unchanged or changed by this explanation, and it compares every output in the same way.

\textbf{The explanation, in full.} The 7+1 spacetime has one time-like dimension ($x_0$, the direction of the octonion unit $e_0=1$) and seven space-like dimensions ($x_1,x_2,x_3$ ordinary space, $x_4,\dots,x_7$ auxiliary space). Two different objects describe this, and neither may be replaced by the other. (i) The tensor package xAct, used in section 12 of 7+1.nb exactly as in the author’s notebook Einstein-Lovelock-7+1-Nash.nb, declares the metric with \texttt{DefMetric[\{7, 1, 0\}, g71[-a, -b], CD, …]}: the list $\{7,1,0\}$ declares seven space-like dimensions, one time-like dimension and no degenerate dimension — the numbers of entries $+1$ and $-1$ of $\eta_{71}=\mathrm{diag}(-1,+1,\dots,+1)$ (\texttt{η71}), the same spacetime written in the sign convention of 4+4.nb (space-like $=+1$, time-like $=-1$). (ii) The matrix $\eta=-\eta_{71}=\mathrm{diag}(+1,-1,\dots,-1)$ (\texttt{η17}) is simply the squares of the octonion basis: $e_0^2=+1$ for the unit $e_0=1$ (time) and $e_i^2=-1$ for the seven imaginary units $e_1,\dots,e_7$ (space); more generally it is the real part of the product of two basis units, and it gives the quadratic form of the octonions:
\begin{gather*}
\eta_{AB}=\mathrm{Re}(e_Ae_B),\qquad e_0^2=+1,\qquad e_i^2=-1\ \ (i=1,\dots,7),\\ \mathrm{Re}(x\,x)=x_0^2-x_1^2-\cdots-x_7^2=\eta_{AB}\,x^Ax^B\qquad (x=x^Ae_A\in\mathbb O).
\end{gather*}
$\eta$ gives the signs of the Clifford relation $\Gamma_A\Gamma_B+\Gamma_B\Gamma_A=2\,\eta_{AB}\,1_{16}$ of the real 16×16 Dirac matrices: time squares to $+1_{16}$ and each of the seven space directions to $-1_{16}$, exactly like the octonion units (in the notation $Cl(p,q)$, with $p$ generators squaring to $+1$ and $q$ to $-1$, the algebra $Cl(1,7)\cong\mathbb R(16)$). It also gives the signs of the metric components: $g_{00}>0$ along the time-like direction, negative squares along the space-like directions. Every component computed by 7+1.nb follows from the component matrix of the metric; the declared list fixes for xAct the sign of the determinant of the metric, \texttt{SignDetOfMetric[g71]} $=-1$, which is the same in both sign conventions, since $\det \eta=\det \eta_{71}=-1$. 7+1.nb checks \texttt{SignatureOfMetric[g71]} $=\{7,1,0\}$, \texttt{SignDetOfMetric[g71]} $=-1$, $\mathrm{Re}(e_Ae_B)=\eta_{AB}$ and $\mathrm{Re}(x\,x)=\eta_{AB}x^Ax^B$.

\Needspace{12\baselineskip}\section*{2. The author’s note}\addcontentsline{toc}{section}{2. The author’s note}
The note is a PDF printout of an e-mail of 28 October 2026 (07:15) from the author to himself, collecting a discussion of the radion. It is the author’s private material: it is not part of the repository (the documents paraphrase it and quote only its formulas). Its content, in the order of the note:

\begin{enumerate}
\item the Kaluza–Klein ansatz $d\hat s^2=e^{2\alpha\phi}g_{\mu\nu}dx^\mu dx^\nu+e^{2\beta\phi}\gamma_{ij}dy^idy^j$ for $D=4+n$ (signature $(-,+,+,+)$), the Einstein-frame values $\alpha=-\sqrt{n/(2(n+2))}$, $\beta=-\frac2n\alpha=\sqrt{2/(n(n+2))}$, the reduced action $\frac{V_n}{2\kappa_D^2}\int\sqrt{-g}[R-\frac12(\partial\phi)^2]$ and the physical volume $V_{\rm phys}=V_ne^{n\beta\phi}$; the radion as a modulus, the fifth-force problem, stabilization;
\item the Randall–Sundrum model RS1: $ds^2=e^{-2kT(x)|y|}g-T(x)^2dy^2$, $S_5=M_5^3\int d^4x\int_{-\pi}^{\pi}dy\sqrt{-G}R_5$, $\mu=f_\mu e^{-k\pi T}$ with $f_\mu=\sqrt{24M_5^3/k}$, $M_{\rm Pl}^2=(M_5^3/k)(1-e^{-2k\pi\langle T\rangle})$, the TeV-scale vacuum value and coupling;
\item an eight-dimensional explanation for students with $n=4$: the values $\alpha=-1/\sqrt3$ and $\beta=1/\sqrt{12}$, the action $S_{\rm eff}=\int\sqrt{-g}[\frac{M_{\rm Pl}^2}2R-\frac12(\partial\phi)^2]$, and the volume $V_{\rm phys}=V_{\rm baseline}e^{(2/\sqrt3)\phi}$;
\item a Goldberger–Wise stabilization in eight dimensions: the potential $V(\phi)=\int d^4y\sqrt\gamma\,e^{4\beta\phi}\,u$ with the density $u=\frac12e^{-2\beta\phi}(\nabla_y\Phi)^2+\frac12m_8^2\Phi^2$, the mass $m_{\rm radion}^2=V''(\phi_0)$, the estimate $m_{\rm radion}\sim(m_8/M_8)^2M_{4D}$ and the Yukawa force $F\propto e^{-m_{\rm radion}r}/r^2$;
\item a warped eight-dimensional metric $ds^2=e^{-2A(y)T(x)}g-T(x)^2\gamma_{ij}dy^idy^j$ with the equations $3A''+6A'^2=-\kappa_8^2\Lambda_8$, $6A'^2=-\kappa_8^2\Lambda_8$, the solution $A=k|y|$, the tensions $\lambda_A=-\lambda_B=6k/\kappa_8^2$;
\item the Israel junction conditions $[K_{\mu\nu}-\hat g_{\mu\nu}K]^+_-=\kappa_8^2S_{\mu\nu}$, $S_{\mu\nu}=-\lambda\hat g_{\mu\nu}$, $K^\pm_{\mu\nu}=\mp k\hat g_{\mu\nu}$, ending with $\lambda_{\rm UV}=-6k/\kappa_8^2$, $\lambda_{\rm IR}=+6k/\kappa_8^2$;
\item the radion as the distance between the branes: $S_8=\frac1{2\kappa_8^2}\int d^4x\int_0^\pi dy\sqrt{-\hat g}\hat R_8$, $S_{\rm eff}=\frac1{2\kappa_8^2}\int\sqrt{-g}[\frac1k(1-e^{-2k\pi T})R_4+\frac6k\partial(e^{-k\pi T})\partial(e^{-k\pi T})]$, $f_\mu=\sqrt{12/(\kappa_8^2k)}$, $M_{\rm Pl}^2=(1-e^{-2k\pi\langle T\rangle})/(\kappa_8^2k)$;
\item the Goldberger–Wise mechanism in the warped model: $\Phi''-4k\Phi'-m^2\Phi=0$, $\Phi=e^{2ky}(C_1e^{\nu ky}+C_2e^{-\nu ky})$, $\nu=\sqrt{4+m^2/k^2}\approx2+\epsilon$, $\epsilon=m^2/(4k^2)$, $V(T)\approx k\epsilon v_A^2+4ke^{-4k\pi T}(v_B-v_Ae^{-\epsilon k\pi T})^2$, $k\pi\langle T\rangle\approx(4k^2/m^2)\ln(v_A/v_B)$, and the radion mass $m_{\rm radion}^2=\frac23\kappa_8^2k^2\epsilon^2v_B^2(\langle\mu\rangle/f_\mu)^2$;
\item Higgs–radion mixing: $S_{\rm mixing}=-\xi\int\sqrt{-g_{\rm ind}}R_4(g_{\rm ind})H^\dagger H$, $\mathcal L_{\rm kinetic}=-\frac12(\partial h_0)^2-\frac12(\partial r_0)^2-6\xi\gamma\,\partial h_0\cdot\partial r_0$, $\gamma=v/\langle\mu\rangle$, $Z^2=1+6\xi\gamma^2(1-6\xi)$, the shift $h_0=h'-\frac{6\xi\gamma}Zr'$, $r_0=r'/Z$, the rotation by $\theta$ with $\tan2\theta=12\gamma\xi Zm_{h_0}^2/(m_{\rm radion}^2-m_{h_0}^2(Z^2-36\xi^2\gamma^2))$;
\item the trace-anomaly couplings ($b_{\rm EM}=-11/3$, $b_{\rm QCD}=7$), $C_h=\cos\theta-\frac{6\xi\gamma}Z\sin\theta$, $C_r=\sin\theta/Z$, the rate ratios $|C_h+C_r\gamma b/F|^2$ for $h\to\gamma\gamma$ ($F_{\rm SM}\approx-6.5$) and $h\to gg$ ($F_{\rm heavy}\approx-4/3$), and LHC bounds ($\langle\mu\rangle>2.5$–$3$ TeV; $-0.1<\xi<0.4$ for $\langle\mu\rangle\approx3$ TeV, $\gamma\approx0.08$).
\end{enumerate}
The note’s file name is excluded from the repository together with the other private files; the copy used was read on 28 October 2026 at 07:23 (text extracted with pdftotext).

\Needspace{12\baselineskip}\section*{3. Decisions}\addcontentsline{toc}{section}{3. Decisions}
{\footnotesize
\begin{longtable}{>{\raggedright\arraybackslash}p{\dimexpr 0.0500\linewidth-2\tabcolsep\relax}>{\raggedright\arraybackslash}p{\dimexpr 0.4700\linewidth-2\tabcolsep\relax}>{\raggedright\arraybackslash}p{\dimexpr 0.4800\linewidth-2\tabcolsep\relax}}
\toprule
\textbf{\#} & \textbf{decision} & \textbf{reason / equation} \\
\midrule\endfirsthead
\toprule
\textbf{\#} & \textbf{decision} & \textbf{reason / equation} \\
\midrule\endhead
\endfoot
\bottomrule\endlastfoot
1 & append section 26 to 7+1.nb ("The radion: the size of the extra space, on the 4-torus and on the 4-sphere"); the summary of all checks becomes section 27 & every worked example and every check of the note is then computed by the notebook itself, in the same evaluation as everything else \\
2 & write an independent Python program, tools/ref\_radion.py & its own curvature routine (sympy), Euler–Lagrange operators, lattice sieve (numpy) and 60-digit arithmetic (mpmath); no code shared with the notebook \\
3 & back up the evaluated notebook before appending section 26 and check that sections 0–25 are unchanged & tools/verify\_7p1\_nb.py compares every cell and output with the backup (section 7.3) \\
4 & the radion ansatz in the Einstein frame of 7+1.nb: $ds^2=e^{-4\beta}g_E-e^{2\beta}\gamma$ & the conventions of section 20.3 ($\beta=\ln(b/b_0)$, $\phi=\sqrt{12}M_{\rm Pl}\beta$); the note’s $\beta\phi$ is our $\beta$, its $\alpha\phi$ our $-2\beta$ \\
5 & the 4-sphere: radius $r=r_0e^\beta$, angles $x_4,x_5,x_6\in[0,\pi]$, $x_7\in[0,2\pi)$, the brane at the north pole $x_4=0$ & the standard hyperspherical coordinates; the north pole is the point fixed by $SO(4)$ \\
6 & the reduction is checked on homogeneous metrics with arbitrary $N(x_0)$, $a(x_0)$, $\beta(x_0)$: the Euler–Lagrange expressions of the difference of the two actions must vanish identically & a difference of Lagrangians that is a total derivative has vanishing Euler–Lagrange expressions; covariance extends the result to $\beta(x)$ \\
7 & printed numbers are computed (solved coefficients, solutions of the conditions), never typed & the outputs of section 26 are compared with those of the Python program (section 7.2) \\
8 & the 4-form flux of the sphere example is an illustration & 7+1.nb has no such field; it shows how a sphere can be stabilized \\
9 & the warped example is the five-dimensional RS1, in the conventions $G^A{}_B-\Lambda\delta^A_B=\kappa T^A{}_B$ & the note’s warped formulas are five-dimensional; the eight-dimensional statement is tested separately (no solution) \\
10 & tensions from the $\delta$-function part of the field equations, and the junction condition checked against it & this form involves no choice of normal; it decides the sign of the junction condition \\
11 & the Higgs–radion mixing in the conventions of Giudice, Rattazzi and Wells, translated to the signature $(+,-,-,-)$ & the note follows them; in this signature the curvature term of the brane Lagrangian is $+\frac\xi2R(h)\varphi_H^2$ \\
12 & the numbers of the mixing in 7+1 in 60-digit arithmetic & the rate changes, $\sim10^{-31}$, are below double precision \\
13 & the numbers today: the sphere is given the volume $V_\perp=M_{\rm Pl}^2/M_8^6$ of the torus & the same four-dimensional Newton constant, so that the two spaces can be compared \\
14 & the trace-anomaly coefficients and the LHC bounds are quoted from the note and not refitted & they concern the warped radion; for 7+1 the relevant parameter $\gamma=v/(\sqrt3M_{\rm Pl})$ is computed \\*
\end{longtable}}
\Needspace{12\baselineskip}\section*{4. Inputs from the earlier sections of 7+1.nb}\addcontentsline{toc}{section}{4. Inputs from the earlier sections of 7+1.nb}
{\footnotesize
\begin{longtable}{>{\raggedright\arraybackslash}p{\dimexpr 0.2600\linewidth-2\tabcolsep\relax}>{\raggedright\arraybackslash}p{\dimexpr 0.2200\linewidth-2\tabcolsep\relax}>{\raggedright\arraybackslash}p{\dimexpr 0.1000\linewidth-2\tabcolsep\relax}>{\raggedright\arraybackslash}p{\dimexpr 0.4200\linewidth-2\tabcolsep\relax}}
\toprule
\textbf{object} & \textbf{Mathematica} & \textbf{section} & \textbf{equation or value} \\
\midrule\endfirsthead
\toprule
\textbf{object} & \textbf{Mathematica} & \textbf{section} & \textbf{equation or value} \\
\midrule\endhead
\endfoot
\bottomrule\endlastfoot
curvature routine, cosmological metric & \texttt{curvature[]}, \texttt{gCos}, \texttt{cvCos} & 11 & $ds^2=dx_0^2-a^2d\vec x^2-b^2d\vec y^2$ and its $R$, $G^A{}_B$ \\
rewriting in expansion rates & \texttt{toH} & 11 & $a'\to H_aa$, $a''\to(H_a'+H_a^2)a$, the same for $b$ \\
brane spinor matrices & \texttt{rhoK} & 17.2 & $\rho_k\rho_l+\rho_l\rho_k=-2\delta_{kl}$ \\
electroweak constants & \texttt{vEW}, \texttt{gFermi}, \texttt{mHb} & 17.3 & $v=(\sqrt2G_F)^{-1/2}$, $m_h=125.20$ GeV \\
stress tensor of a comoving energy & \texttt{feStress} & 20.2 & $\varepsilon=W/(a^3b^4)$, $p_a=-\partial_aW/(3a^2b^4)$, $p_b=-\partial_bW/(4a^3b^3)$ \\
radion, its equation, mass and coupling & — & 20.3 & $\beta''+3H_a\beta'+4\beta'^2=\kappa(\frac\varepsilon6-\frac{p_a}2+\frac{p_b}3)+\frac\Lambda3$, $m_r^2=\kappa K_U/12$, $\phi T^\mu{}_\mu/(\sqrt3M_{\rm Pl})$ \\
Kasner exponents, capture, decay, the radion today and in the far future & — & 21–23 & the values quoted in section 10 of the document \\
Planck masses, volume, lattice & \texttt{rIIMPl}, \texttt{rIIM8}, \texttt{rIIVp}, \texttt{rIIr4} & 22.2 & $M_{\rm Pl}=2.435\times10^{18}$ GeV, $M_8=10^4$ GeV, $V_\perp=M_{\rm Pl}^2/M_8^6$ \\
the torus today, $\hbar c$ & \texttt{tbLp}, \texttt{tbHbarCm}, \texttt{tbNmu} & 25 & $\ell_\perp=V_\perp^{1/4}$, $\hbar c=1.973269804\times10^{-16}$ GeV m \\*
\end{longtable}}
Before section 26 was appended, the evaluated notebook (sections 0–25 and the summary of all checks, then numbered 26) was copied to \texttt{backups/\allowbreak{}7+1\_\allowbreak{}2026-10-28\_\allowbreak{}before\_\allowbreak{}radion.nb} (2031601 bytes, SHA-256 \texttt{44385a6b\allowbreak{}1d808f1f\allowbreak{}95b1d634\allowbreak{}d29c3960\allowbreak{}c6251cab\allowbreak{}96f436b9\allowbreak{}927522ab\allowbreak{}e0fd8397}; excluded from the repository by .gitignore).

\Needspace{12\baselineskip}\section*{5. Data and their sources}\addcontentsline{toc}{section}{5. Data and their sources}
{\footnotesize
\begin{longtable}{>{\raggedright\arraybackslash}p{\dimexpr 0.2500\linewidth-2\tabcolsep\relax}>{\raggedright\arraybackslash}p{\dimexpr 0.4000\linewidth-2\tabcolsep\relax}>{\raggedright\arraybackslash}p{\dimexpr 0.3500\linewidth-2\tabcolsep\relax}}
\toprule
\textbf{quantity} & \textbf{value used} & \textbf{source} \\
\midrule\endfirsthead
\toprule
\textbf{quantity} & \textbf{value used} & \textbf{source} \\
\midrule\endhead
\endfoot
\bottomrule\endlastfoot
Eöt-Wash bound & $\lambda<38.6\,\mu$m for $|\alpha|=1$ (95\% CL) & J. G. Lee et al., Phys. Rev. Lett. 124 (2020) 101101 \\
$\hbar c$ & $1.973269804\times10^{-16}$ GeV m & CODATA 2018 \\
$M_{\rm Pl}$ (reduced), $M_8$, $m_r$ & $2.435\times10^{18}$ GeV, 10 TeV, 10 GeV & as in sections 17–25 \\
$G_F$, $m_h$ & $1.1663788\times10^{-5}$ GeV$^{-2}$, 125.20 GeV & Particle Data Group (as in section 17.3) \\
$b_{\rm EM}$, $F_{\rm SM}$, $b_{\rm QCD}$, $F_{\rm heavy}$ & $-11/3$, $-6.5$, $7$, $-4/3$ & the note (standard values of the Standard Model) \\
$\langle\mu\rangle$ of the warped example & 3 TeV & the note \\
RS1 conventions & $2M^3\int\sqrt{-G}R$, $V_{\rm hid}=-V_{\rm vis}=24M^3k$, $\Lambda=-24M^3k^2$ & L. Randall and R. Sundrum, Phys. Rev. Lett. 83 (1999) 3370 \\
Goldberger–Wise results & $f=\sqrt{24M^3/k}$, $m_\varphi^2=k^2v_v^2\epsilon^2e^{-2kr_c\pi}/(3M^3)$ & W. D. Goldberger and M. B. Wise, Phys. Lett. B 475 (2000) 275 \\
Dirac spectrum of $S^4$ & $\pm(l+2)/r$, degeneracy $4\binom{l+3}3$ & R. Camporesi and A. Higuchi, J. Geom. Phys. 20 (1996) 1 \\*
\end{longtable}}
Software: Wolfram Language 15.0.1 (front end WolframNB.exe; the kernel configuration "Eval71" is a fresh kernel); Python 3.13.17 with numpy 2.5.3, sympy 1.14.0, mpmath 1.3.0.

\Needspace{12\baselineskip}\section*{6. Derivations}\addcontentsline{toc}{section}{6. Derivations}
\textbf{Einstein frame.} Under $g\to\Omega^2g$ in four dimensions $\sqrt{|g|}\to\Omega^4\sqrt{|g|}$ and $R\to\Omega^{-2}R$. For the blocks $e^{2\alpha}g$ and $e^{2\beta}\gamma$ ($n$ extra dimensions) the term $\sqrt{|G|}R[G]$ contains $e^{4\alpha}e^{n\beta}e^{-2\alpha}\sqrt{|g|}\sqrt\gamma R[g]$, constant iff $2\alpha+n\beta=0$; for $n=4$: $\alpha=-2\beta$, the factor $e^{-4\beta}$.

\textbf{The kinetic coefficient.} With the ansatz $G=\mathrm{diag}(N^2e^{-n\beta},-a^2e^{-n\beta}\cdot1_3,-e^{2\beta}\cdot1_n)$ (per unit coordinate volume of $T^n$), the reduction has the form
\begin{equation*}
-\frac1{2\kappa}\sqrt{|G|}R=-\frac1{2\kappa}Na^3R_E+\frac1{2\kappa}\,c_n\,\frac{a^3\beta'^2}N+\text{(total derivative)},
\end{equation*}
and the coefficient is SOLVED from the Euler–Lagrange expression for $N$ (it enters linearly): $c_n=n(n+2)/2$, i.e. $\frac32$, $4$, $\frac{15}2$, $12$ for $n=1,\dots,4$. With $\beta=\beta_n\varphi$ the kinetic term is $\frac{M_{\rm Pl}^2}2c_n\beta_n^2(\partial\varphi)^2$: the note’s $\beta_n^2=2/(n(n+2))$ gives $c_n\beta_n^2=1$, the form with $\frac12$ needs $\beta_n^2=1/(n(n+2))$. For $n=4$: $6M_{\rm Pl}^2(\partial\beta)^2=\frac12(\partial\phi)^2$ with $\phi=\sqrt{12}M_{\rm Pl}\beta$.

\textbf{The potential.} $\sqrt{|G|}=e^{-4\beta}\sqrt{|g_E|}\sqrt\gamma$, and the internal block $-e^{2\beta}\gamma$ contributes $-e^{-2\beta}R_{\rm int}[\gamma]$ to $R$ in the signature of 7+1.nb. Hence $-\frac1{2\kappa}\int\sqrt{|G|}(R+2\Lambda)-\int\sqrt{|G|}U\supset\sqrt{|g_E|}e^{-4\beta}[\frac{V_0}{2\kappa}R_{\rm int}(\beta)-M_{\rm Pl}^2\Lambda-V_0U]=-\sqrt{|g_E|}V(\beta)$ with
\begin{equation*}
V(\beta)=e^{-4\beta}\Big[M_{\rm Pl}^2\Lambda+V_0U-\frac{M_{\rm Pl}^2}2R_{\rm int}\Big],\qquad R_{\rm int}(S^4)=\frac{12}{r_0^2e^{2\beta}},\qquad m_r^2=\frac{V''(0)}{12M_{\rm Pl}^2}.
\end{equation*}
\textbf{The four cases.} With $U=U_0+U_1\beta+\frac{K_U}2\beta^2$ and $V=M_{\rm Pl}^2[e^{-4\beta}(\Lambda+\kappa U)-6\sigma e^{-6\beta}/r_0^2+Fe^{-12\beta}]$, the conditions $V(0)=V'(0)=0$ and $V''(0)$ give: $T^4$ ($\sigma=F=0$): $\Lambda+\kappa U_0=0$, $U_1=0$, $m_r^2=\kappa K_U/12$; $S^4$ with $U$: $(\Lambda+\kappa U_0)r_0^2=6$, $\kappa U_1r_0^2=-12$, $V''(0)=M_{\rm Pl}^2(\kappa K_U-24/r_0^2)$, $m_r^2=\kappa K_U/12-2/r_0^2$; $S^4$ with $\Lambda$ only: $V'(0)=0\Rightarrow\Lambda r_0^2=9$, $V(0)=3M_{\rm Pl}^2/r_0^2$, $H^2=1/r_0^2$, $V''(0)=-72M_{\rm Pl}^2/r_0^2$, $m^2=-6H^2$; $S^4$ with flux: $\Lambda-6/r_0^2+F=0$ and $-4\Lambda+36/r_0^2-12F=0$ give $Fr_0^2=\frac32$, $\Lambda r_0^2=\frac92$, $V''(0)=72M_{\rm Pl}^2/r_0^2$, $m^2=6/r_0^2$. The eight-dimensional radion equation $\beta''+3H_a\beta'+4\beta'^2=\kappa(\frac\varepsilon6-\frac{p_a}2+\frac{p_b}3)+\frac\Lambda3-\frac{3\sigma}{b^2}$ reproduces each: its right-hand side vanishes at $\beta=0$ and its derivative is $-m^2$.

\textbf{The coupling and the fifth force.} On the brane $h=e^{-4\beta}g_E$, so $-M\int d\tau=-M\int e^{-2\beta}dx_0$ and the linear term is $2M\beta=M\phi/(\sqrt3M_{\rm Pl})$. A scalar with coupling $g$ gives $-g^2MM'/(4\pi r)=-2g^2M_{\rm Pl}^2\,G_4MM'/r$ ($8\pi G_4=M_{\rm Pl}^{-2}$), so $\alpha=2(gM_{\rm Pl})^2=\frac23$. The Eöt-Wash bound gives $m_r>\hbar c/38.6\,\mu{\rm m}=1.973269804\times10^{-7}\,{\rm eV\,m}/3.86\times10^{-5}\,{\rm m}$.

\textbf{Levels.} On $S^4$, $N=\sum_{l\le L}d_l=(L+1)(L+2)^2(L+3)/12=y^2(y^2-1)/12$ with $y=L+2$; for integer $x=\mu r>2$ the largest $l$ with $l(l+3)\le x^2$ is $L=x-2$ (since $(x-2)(x+1)\le x^2<(x-1)(x+2)$), so $N/(x^4/12)-1=-1/x^2$. The zonal harmonic: $\int_{S^4}(C_l/C_l(1))^2=V/d_l$ (checked for $l\le6$), so $Z_l(0)=\sqrt{d_l/V}$.

\textbf{Killing spinors.} With $\tilde\gamma_a=i\rho_{a+3}$ Hermitian and the frame $e^a$, the spinor $\varepsilon_s$ of the document satisfies $\nabla_i\varepsilon_s=s\frac i{2r}e_i{}^a\tilde\gamma_a\varepsilon_s$ identically (checked exactly after writing all angles as twice half-angles); then $\not D\varepsilon_s=s\frac{2i}r\varepsilon_s$, saturating $|\lambda|^2\ge nR_{\rm int}/(4(n-1))=4/r^2$.

\textbf{RS1.} For $ds^2=e^{-2A}\eta-dy^2$: $G^\mu{}_\nu=3A''-6A'^2$, $G^y{}_y=-6A'^2$. $A=k|y|$: $A''=2k\delta(y)$ at the UV brane; the $\delta$ part $6k\delta(y)\delta^\mu_\nu=\kappa f\delta(y)\delta^\mu_\nu$ gives $f_{\rm UV}=6k/\kappa$. With $K^\mu{}_\nu=\frac12g^{\mu\alpha}\partial_yg_{\alpha\nu}=-A'\delta^\mu_\nu$: $[K^\mu{}_\nu]-\delta^\mu_\nu[K]=-[A']+4[A']=3[A']=6k$, equal to the $\delta$ part of $G^\mu{}_\nu$; in the signature $(-,+,+,+,+)$ $G^A{}_B$ changes sign, so there $[K_{\mu\nu}-\hat g_{\mu\nu}K]=-\kappa S_{\mu\nu}$. $M_{\rm Pl}^2=\frac1\kappa\int_{-\pi R}^{\pi R}e^{-2k|y|}dy$. For $ds^2=e^{-2kT|y|}g-T^2dy^2$ the bulk Lagrangian, integrated over $y$, is $c_2(T)T''+\dots$; subtracting the total derivative $(c_2T')'$ leaves $\frac3{\kappa k}(\partial e^{-k\pi T})^2$, and $\frac12f^2=\frac3{\kappa k}$.

\textbf{Goldberger–Wise.} $L=e^{-4ky}(-\frac12\Phi'^2-\frac12m^2\Phi^2)$ gives $e^{-4ky}(\Phi''-4k\Phi'-m^2\Phi)=0$; the roots of $\lambda^2-4k\lambda-m^2=0$ are $k(2\pm\nu)$. For $V=\frac{4k}{f^4}\mu^4W^2$, $W=v_B-v_A(\mu/f)^\epsilon$: $V'=\frac{4k}{f^4}(4\mu^3W^2+2\mu^4WW')$ vanishes at $W=0$; there $V''=\frac{8k}{f^4}\mu^4W'^2$ with $W'=-\epsilon v_B/\mu$, so $V''=8k\epsilon^2v_B^2\mu^2/f^4$, and with $f^2=6/(\kappa k)$: $\frac43\kappa k^2\epsilon^2v_B^2(\mu/f)^2$.

\textbf{Higgs–radion.} For $h=\Omega^2\eta$ in four dimensions, $\sqrt{|h|}R(h)=-6\Omega\Box\Omega$ (computed). With $\Omega=1+o_1r+o_2r^2$ and $\varphi_H=v+h$, the curvature term $\frac\xi2(-6\Omega\Box\Omega)\varphi_H^2$ has the quadratic part $-3\xi(v^2o_1^2r\Box r+2vo_1h\Box r)$ plus total derivatives ($o_2$ enters only as $\Box(r^2)$); integrating by parts gives $3\xi\gamma^2(\partial r)^2+6\xi\gamma\,\partial h\cdot\partial r$ ($\gamma=vo_1$). $S^{\mathsf T}KS=1$ requires $(K_{rr}-36\xi^2\gamma^2)/Z^2=1$. In $Q^{\mathsf T}M'Q$ the off-diagonal element is $-\frac12(M'_{11}-M'_{22})\sin2\theta+M'_{12}\cos2\theta$.

\Needspace{12\baselineskip}\section*{7. Verification}\addcontentsline{toc}{section}{7. Verification}
\Needspace{8\baselineskip}\subsection*{7.1 The checks of section 26 of 7+1.nb}\addcontentsline{toc}{subsection}{7.1 The checks of section 26 of 7+1.nb}
\textbf{Evaluation of 7+1.nb}: Wolfram 15.0.1 for Microsoft Windows (64-bit) (July 2, 2026), a fresh kernel (front-end kernel configuration Eval71); finished Wed 28 Oct 2026 08:57:47; 283.883 s; message cells 0; output cells 554. The checks of section 26:

{\footnotesize
\begin{longtable}{>{\raggedright\arraybackslash}p{\dimexpr 0.0550\linewidth-2\tabcolsep\relax}>{\raggedright\arraybackslash}p{\dimexpr 0.7450\linewidth-2\tabcolsep\relax}>{\raggedright\arraybackslash}p{\dimexpr 0.0800\linewidth-2\tabcolsep\relax}>{\raggedright\arraybackslash}p{\dimexpr 0.1200\linewidth-2\tabcolsep\relax}}
\toprule
\textbf{\#} & \textbf{check} & \textbf{result} & \textbf{section} \\
\midrule\endfirsthead
\toprule
\textbf{\#} & \textbf{check} & \textbf{result} & \textbf{section} \\
\midrule\endhead
\endfoot
\bottomrule\endlastfoot
1 & \texttt{section 26 engine for diagonal metrics: rdScal and rdEin reproduce curvature[] of section 11 for the cosmological metric (R and every G\textasciicircum{}A\_\allowbreak{}A)} & True & 26.1 \\
2 & \texttt{the 4-sphere of radius r0: Ricci\_\allowbreak{}ij = 3 gamma\_\allowbreak{}ij (gamma the unit-sphere metric), Ricci\_\allowbreak{}ij = 0 for i != j, so R(S\textasciicircum{}4) = 12/\allowbreak{}r0\textasciicircum{}2 (in the signature of 7+1.nb the space-like block gives R = -12/\allowbreak{}r0\textasciicircum{}2, also with curvature[] of section 11); the volume of the unit S\textasciicircum{}4 is 8 Pi\textasciicircum{}2/\allowbreak{}3} & True & 26.1 \\
3 & \texttt{radion on T\textasciicircum{}4: for ds\textasciicircum{}2 = e\textasciicircum{}(-4 beta) (N\textasciicircum{}2 dx0\textasciicircum{}2 - a\textasciicircum{}2 dx.dx) - e\textasciicircum{}(2 beta) b0\textasciicircum{}2 dy.dy the reduced action -(1/\allowbreak{}(2 kappa)) Int(R + 2 Lambda) - Int U equals -(M\_\allowbreak{}Pl\textasciicircum{}2/\allowbreak{}2) R\_\allowbreak{}E + 6 M\_\allowbreak{}Pl\textasciicircum{}2 (d beta)\textasciicircum{}2 - e\textasciicircum{}(-4 beta)(M\_\allowbreak{}Pl\textasciicircum{}2 Lambda + V0 U) up to a total derivative (Euler-Lagrange expressions for N, a, beta vanish), M\_\allowbreak{}Pl\textasciicircum{}2 = V0/\allowbreak{}kappa, V0 = b0\textasciicircum{}4 Lc\textasciicircum{}4} & True & 26.1 \\
4 & \texttt{radion on S\textasciicircum{}4: the same reduction for ds\textasciicircum{}2 = e\textasciicircum{}(-4 beta) g\_\allowbreak{}E - e\textasciicircum{}(2 beta) r0\textasciicircum{}2 gamma (unit-sphere gamma) holds with the additional potential -(M\_\allowbreak{}Pl\textasciicircum{}2/\allowbreak{}2) e\textasciicircum{}(-4 beta) 12/\allowbreak{}(r0\textasciicircum{}2 e\textasciicircum{}(2 beta)) = -6 M\_\allowbreak{}Pl\textasciicircum{}2 e\textasciicircum{}(-6 beta)/\allowbreak{}r0\textasciicircum{}2 from the curvature of the sphere, V0 = (8 Pi\textasciicircum{}2/\allowbreak{}3) r0\textasciicircum{}4; R is independent of the angles} & True & 26.1 \\
5 & \texttt{n extra dimensions (torus T\textasciicircum{}n, D = 4 + n, n = 1..4, Einstein frame e\textasciicircum{}(-n B) g\_\allowbreak{}E): the kinetic term is (M\_\allowbreak{}Pl\textasciicircum{}2/\allowbreak{}2)(n(n+2)/\allowbreak{}2)(dB)\textasciicircum{}2; with the note's B = beta\_\allowbreak{}n phi, beta\_\allowbreak{}n = Sqrt[2/\allowbreak{}(n(n+2))], it is (M\_\allowbreak{}Pl\textasciicircum{}2/\allowbreak{}2)(d phi)\textasciicircum{}2, not (M\_\allowbreak{}Pl\textasciicircum{}2/\allowbreak{}2)(1/\allowbreak{}2)(d phi)\textasciicircum{}2; the form with 1/\allowbreak{}2 needs beta\_\allowbreak{}n = 1/\allowbreak{}Sqrt[n(n+2)]} & True & 26.1 \\
6 & \texttt{the note's eight-dimensional values alpha = -1/\allowbreak{}Sqrt[3], beta = 1/\allowbreak{}Sqrt[12]: 2 alpha + 4 beta == 0 (Einstein frame), 4 beta == 2/\allowbreak{}Sqrt[3] (V = V0 e\textasciicircum{}(4 beta phi)), and Sqrt[12] beta == 1, so M\_\allowbreak{}Pl phi\_\allowbreak{}note = Sqrt[12] M\_\allowbreak{}Pl ln(b/\allowbreak{}b0) is the canonical radion of section 20.3} & True & 26.1 \\
7 & \texttt{the cosmological metric with the sphere of radius b(x0) (ds\textasciicircum{}2 = dx0\textasciicircum{}2 - a\textasciicircum{}2 dx.dx - b\textasciicircum{}2 dOmega\_\allowbreak{}4\textasciicircum{}2): G\textasciicircum{}A\_\allowbreak{}A = (torus value) + (6, 6, 6, 6, 3, 3, 3, 3)/\allowbreak{}b\textasciicircum{}2 and R = (torus value) - 12/\allowbreak{}b\textasciicircum{}2} & True & 26.2 \\
8 & \texttt{radion equation on S\textasciicircum{}4 (Einstein gravity; 11 and 44 equations with the 00 constraint): Hb' + 3 Ha Hb + 4 Hb\textasciicircum{}2 == kappa (eps/\allowbreak{}6 - pa/\allowbreak{}2 + pb/\allowbreak{}3) + Lambda/\allowbreak{}3 - 3/\allowbreak{}b\textasciicircum{}2} & True & 26.2 \\
9 & \texttt{radion on T\textasciicircum{}4 with U = U0 + U1 beta + (K\_\allowbreak{}U/\allowbreak{}2) beta\textasciicircum{}2: a flat vacuum at b = b0 needs Lambda + kappa U0 = 0 and U1 = 0, and m\_\allowbreak{}r\textasciicircum{}2 = V''(beta)/\allowbreak{}(12 M\_\allowbreak{}Pl\textasciicircum{}2) = kappa K\_\allowbreak{}U/\allowbreak{}12 (section 20.3)} & True & 26.2 \\
10 & \texttt{radion on S\textasciicircum{}4 with U: Lambda + kappa U0 == 6/\allowbreak{}r0\textasciicircum{}2, kappa U1 == -12/\allowbreak{}r0\textasciicircum{}2 and m\_\allowbreak{}r\textasciicircum{}2 == kappa K\_\allowbreak{}U/\allowbreak{}12 - 2/\allowbreak{}r0\textasciicircum{}2 from V(beta); the eight-dimensional radion equation gives the same static solution and the same mass} & True & 26.2 \\
11 & \texttt{S\textasciicircum{}4 with Lambda only: dS4 x S\textasciicircum{}4 with r0\textasciicircum{}2 = 9/\allowbreak{}Lambda, H\textasciicircum{}2 = 1/\allowbreak{}r0\textasciicircum{}2 = Lambda/\allowbreak{}9 solves the 00, 11 and 44 equations; it is a MAXIMUM of V: m\textasciicircum{}2 = -6 H\textasciicircum{}2 from V(beta), and m\textasciicircum{}2 = -6/\allowbreak{}r0\textasciicircum{}2 = -6 H\textasciicircum{}2 from the eight-dimensional radion equation} & True & 26.2 \\
12 & \texttt{S\textasciicircum{}4 with a 4-form flux (illustration; W = a\textasciicircum{}3 b\textasciicircum{}4 q\textasciicircum{}2/\allowbreak{}(2 b\textasciicircum{}8) gives eps = rho\_\allowbreak{}F, pa = -rho\_\allowbreak{}F, pb = +rho\_\allowbreak{}F): Minkowski x S\textasciicircum{}4 needs Lambda = 9/\allowbreak{}(2 r0\textasciicircum{}2) and kappa rho\_\allowbreak{}F = 3/\allowbreak{}(2 r0\textasciicircum{}2); m\textasciicircum{}2 = 6/\allowbreak{}r0\textasciicircum{}2 from V(beta) and from the eight-dimensional radion equation} & True & 26.2 \\
13 & \texttt{the note's Goldberger-Wise potential in eight dimensions: for a static profile Phi(y) (here Phi(x4, x5)) Sqrt|G| (1/\allowbreak{}2 G\textasciicircum{}AB dPhi dPhi - 1/\allowbreak{}2 m\textasciicircum{}2 Phi\textasciicircum{}2) == -N a\textasciicircum{}3 e\textasciicircum{}(-4 beta) Sqrt[gamma] (1/\allowbreak{}2 e\textasciicircum{}(-2 beta) |grad Phi|\textasciicircum{}2\_\allowbreak{}gamma + 1/\allowbreak{}2 m\textasciicircum{}2 Phi\textasciicircum{}2): the Einstein-frame factor is e\textasciicircum{}(-4 beta) (e\textasciicircum{}(-4 beta phi) in the note's notation), not e\textasciicircum{}(+4 beta phi)} & True & 26.2 \\
14 & \texttt{the radion couples to a mass M on the brane as (phi/\allowbreak{}(Sqrt[3] M\_\allowbreak{}Pl)) M: -M e\textasciicircum{}(-n beta/\allowbreak{}2) = -M + (n/\allowbreak{}2) M beta + ..., with beta = phi/\allowbreak{}(M\_\allowbreak{}Pl Sqrt[n (n + 2)/\allowbreak{}2]) the coupling is Sqrt[n/\allowbreak{}(2 (n + 2))]/\allowbreak{}M\_\allowbreak{}Pl = 1/\allowbreak{}(Sqrt[3] M\_\allowbreak{}Pl) for n = 4, and the fifth-force strength 2 (coupling M\_\allowbreak{}Pl)\textasciicircum{}2 = n/\allowbreak{}(n + 2) = 2/\allowbreak{}3} & True & 26.3 \\
15 & \texttt{Eot-Wash (|alpha| = 1 excluded for lambda > 38.6 micrometres): a radion of strength 2/\allowbreak{}3 must be heavier than about hbar c/\allowbreak{}38.6 micrometres = 5.1 meV; the 10 GeV radion has the range 2.0 10\textasciicircum{}-17 m; Higgs-radion mixing in 7+1 is governed by gamma = v/\allowbreak{}(Sqrt[3] M\_\allowbreak{}Pl) = 5.8 10\textasciicircum{}-17} & True & 26.3 \\
16 & \texttt{zonal harmonics on S\textasciicircum{}4 (l = 0..6, exact): (1/\allowbreak{}sin\textasciicircum{}3 x4) d/\allowbreak{}dx4 (sin\textasciicircum{}3 x4 d/\allowbreak{}dx4 C\_\allowbreak{}l\textasciicircum{}(3/\allowbreak{}2)(cos x4)) == -l (l + 3) C\_\allowbreak{}l\textasciicircum{}(3/\allowbreak{}2), C\_\allowbreak{}l\textasciicircum{}(3/\allowbreak{}2)(1) == (l + 1)(l + 2)/\allowbreak{}2, and Int over S\textasciicircum{}4 of (C\_\allowbreak{}l/\allowbreak{}C\_\allowbreak{}l(1))\textasciicircum{}2 == V/\allowbreak{}d\_\allowbreak{}l (unit sphere), so the normalized zonal harmonic is Sqrt[d\_\allowbreak{}l/\allowbreak{}V] at the north pole; d\_\allowbreak{}l == Binomial[l + 4, 4] - Binomial[l + 2, 4] for l = 0..20} & True & 26.4 \\
17 & \texttt{Weyl's law on S\textasciicircum{}4: the number of levels with l(l + 3) <= (mu r)\textasciicircum{}2 approaches (mu r)\textasciicircum{}4/\allowbreak{}12 = V mu\textasciicircum{}4/\allowbreak{}(32 Pi\textasciicircum{}2) (V = 8 Pi\textasciicircum{}2 r\textasciicircum{}4/\allowbreak{}3): relative deviation -1.0 10\textasciicircum{}-2, -1.0 10\textasciicircum{}-4, -1.0 10\textasciicircum{}-6 at mu r = 10, 100, 1000 (like -1/\allowbreak{}(mu r)\textasciicircum{}2)} & True & 26.4 \\
18 & \texttt{the brane at the north pole of S\textasciicircum{}4: the embedding (r0 sin x4 sin x5 sin x6 cos x7, ..., r0 cos x4) of S\textasciicircum{}4 in R\textasciicircum{}5 induces r0\textasciicircum{}2 gamma; the north pole x4 = 0 is (0, 0, 0, 0, r0); of the 10 rotations L\_\allowbreak{}ab of R\textasciicircum{}5 (the Killing vectors of S\textasciicircum{}4, SO(5)) the 6 with a, b <= 4 vanish there (SO(4), kept by the brane) and the 4 L\_\allowbreak{}a5 move it} & True & 26.4 \\
19 & \texttt{Killing spinors of S\textasciicircum{}4 (radius r0, frame e\textasciicircum{}a = (r0, r0 sin x4, ...) dx\textasciicircum{}(a+3), gt\_\allowbreak{}a = I rho\_\allowbreak{}(a+3)): eps\_\allowbreak{}s = exp(s (i/\allowbreak{}2) x4 gt\_\allowbreak{}1) exp((1/\allowbreak{}2) x5 gt\_\allowbreak{}1 gt\_\allowbreak{}2) exp((1/\allowbreak{}2) x6 gt\_\allowbreak{}2 gt\_\allowbreak{}3) exp((1/\allowbreak{}2) x7 gt\_\allowbreak{}3 gt\_\allowbreak{}4) eps0 satisfies nabla\_\allowbreak{}mu eps\_\allowbreak{}s == s (i/\allowbreak{}(2 r0)) e\_\allowbreak{}mu\textasciicircum{}a gt\_\allowbreak{}a eps\_\allowbreak{}s (s = +-1, all eps0), so the Dirac operator gives (2 s i/\allowbreak{}r0) eps\_\allowbreak{}s: |lambda| = 2/\allowbreak{}r0, Friedrich's bound n R/\allowbreak{}(4(n - 1)) = 4/\allowbreak{}r0\textasciicircum{}2} & True & 26.4 \\
20 & \texttt{Dirac spectrum on S\textasciicircum{}4: Weyl's law for spinors (4 components) N(|lambda| <= L/\allowbreak{}r) \textasciitilde{} 4 V (L/\allowbreak{}r)\textasciicircum{}4/\allowbreak{}(32 Pi\textasciicircum{}2) = L\textasciicircum{}4/\allowbreak{}3 agrees with the degeneracies 2 x 4 Binomial[l + 3, 3] of the eigenvalues +-i(l + 2)/\allowbreak{}r (to 2.0 10\textasciicircum{}-3 at L = 1000)} & True & 26.4 \\
21 & \texttt{the sphere with the volume of the torus today (V = M\_\allowbreak{}Pl\textasciicircum{}2/\allowbreak{}M8\textasciicircum{}6): r = (3 V/\allowbreak{}(8 Pi\textasciicircum{}2))\textasciicircum{}(1/\allowbreak{}4) = 688.9 GeV\textasciicircum{}-1 = 1.359 10\textasciicircum{}-13 m; lowest scalar levels 4.03 MeV (torus, 8 states) and 2.90 MeV (sphere, l = 1, 5 states); lowest dirac16complex level on the sphere 2/\allowbreak{}r = 2.90 MeV (8 Dirac fermions; the torus has 4 massless ones for m = 0); levels below 5 GeV: torus 1.17335 10\textasciicircum{}13, sphere 1.17375 10\textasciicircum{}13, Weyl 1.17335 10\textasciicircum{}13} & True & 26.4 \\
22 & \texttt{RS1 in five dimensions (mostly minus): G\textasciicircum{}mu\_\allowbreak{}nu = 3 A'' - 6 A'\textasciicircum{}2, G\textasciicircum{}y\_\allowbreak{}y = -6 A'\textasciicircum{}2; bulk A = k y needs Lambda = -6 k\textasciicircum{}2; the delta functions A'' = +-2 k delta at the branes need the tensions f\_\allowbreak{}UV = +6 k/\allowbreak{}kappa, f\_\allowbreak{}IR = -6 k/\allowbreak{}kappa (UV positive)} & True & 26.5 \\
23 & \texttt{Israel junction at the UV brane of RS1: with K\textasciicircum{}mu\_\allowbreak{}nu = (1/\allowbreak{}2) g\textasciicircum{}mu alpha d\_\allowbreak{}y g\_\allowbreak{}alpha nu = -A' delta\textasciicircum{}mu\_\allowbreak{}nu (normal d\_\allowbreak{}y), [K\textasciicircum{}mu\_\allowbreak{}nu] - delta\textasciicircum{}mu\_\allowbreak{}nu [K] = 6 k = (coefficient of A'' in G\textasciicircum{}mu\_\allowbreak{}nu) x (jump of A'), the delta-function part of G\textasciicircum{}mu\_\allowbreak{}nu; in the signature (-,+,+,+,+) of the note G\textasciicircum{}A\_\allowbreak{}B changes sign, so there [K\_\allowbreak{}mu nu - g\_\allowbreak{}mu nu K] = -kappa S\_\allowbreak{}mu nu (the note uses +kappa S\_\allowbreak{}mu nu, hence its reversed tensions)} & True & 26.5 \\
24 & \texttt{RS1: M\_\allowbreak{}Pl\textasciicircum{}2 = (1/\allowbreak{}kappa) Int\_\allowbreak{}(-Pi R)\textasciicircum{}(Pi R) e\textasciicircum{}(-2 k |y|) dy = (1 - e\textasciicircum{}(-2 Pi k R))/\allowbreak{}(kappa k)} & True & 26.5 \\
25 & \texttt{RS1 radion: with R promoted to T(x), the bulk action is (3/\allowbreak{}(kappa k)) (d e\textasciicircum{}(-k Pi T))\textasciicircum{}2 + V(T) + (a total derivative), so the canonical radion is mu = Sqrt[6/\allowbreak{}(kappa k)] e\textasciicircum{}(-k Pi T) (the note: Sqrt[12/\allowbreak{}(kappa k)])} & True & 26.5 \\
26 & \texttt{eight dimensions, warped along one extra direction (ds\textasciicircum{}2 = e\textasciicircum{}(-2 A(y)) eta - dy\textasciicircum{}2 - dz.dz): G\textasciicircum{}mu\_\allowbreak{}nu = 3 A'' - 6 A'\textasciicircum{}2, G\textasciicircum{}y\_\allowbreak{}y = -6 A'\textasciicircum{}2, G\textasciicircum{}z\_\allowbreak{}z = 4 A'' - 10 A'\textasciicircum{}2; with only Lambda (all G\textasciicircum{}A\_\allowbreak{}A equal) A'' = 0 and A' = 0: the note's five-dimensional equations do not carry over} & True & 26.5 \\
27 & \texttt{Goldberger-Wise in RS1: the static bulk scalar (Sqrt|G| = e\textasciicircum{}(-4 k y), G\textasciicircum{}yy = -1) obeys Phi'' - 4 k Phi' - m\textasciicircum{}2 Phi = 0, solved by e\textasciicircum{}(2 k y)(C1 e\textasciicircum{}(nu k y) + C2 e\textasciicircum{}(-nu k y)), nu = Sqrt[4 + m\textasciicircum{}2/\allowbreak{}k\textasciicircum{}2] = 2 + eps + O(m\textasciicircum{}4), eps = m\textasciicircum{}2/\allowbreak{}(4 k\textasciicircum{}2)} & True & 26.5 \\
28 & \texttt{Goldberger-Wise radion mass: V(mu) = 4 k (mu/\allowbreak{}f)\textasciicircum{}4 (v\_\allowbreak{}B - v\_\allowbreak{}A (mu/\allowbreak{}f)\textasciicircum{}eps)\textasciicircum{}2 has V'(<mu>) = 0 at (<mu>/\allowbreak{}f)\textasciicircum{}eps = v\_\allowbreak{}B/\allowbreak{}v\_\allowbreak{}A (k Pi <T> = (4 k\textasciicircum{}2/\allowbreak{}m\textasciicircum{}2) Log[v\_\allowbreak{}A/\allowbreak{}v\_\allowbreak{}B]) and V''(<mu>) = 8 k eps\textasciicircum{}2 v\_\allowbreak{}B\textasciicircum{}2 <mu>\textasciicircum{}2/\allowbreak{}f\textasciicircum{}4; with f\textasciicircum{}2 = 6/\allowbreak{}(kappa k): m\_\allowbreak{}r\textasciicircum{}2 = (4/\allowbreak{}3) kappa k\textasciicircum{}2 eps\textasciicircum{}2 v\_\allowbreak{}B\textasciicircum{}2 (<mu>/\allowbreak{}f)\textasciicircum{}2 (the note's f\textasciicircum{}2 = 12/\allowbreak{}(kappa k) gives 2/\allowbreak{}3)} & True & 26.5 \\
29 & \texttt{conformally flat brane metric h = Omega\textasciicircum{}2 eta (four dimensions, Omega(x0, x1)): Sqrt|h| R(h) == -6 Omega box Omega with box = d0\textasciicircum{}2 - d1\textasciicircum{}2 (curvature conventions of 7+1.nb)} & True & 26.6 \\
30 & \texttt{the Higgs boson on the brane (h = Omega\textasciicircum{}2 eta, Omega = 1 + r/\allowbreak{}Lambda + a2 r\textasciicircum{}2, phi\_\allowbreak{}H = v + h, gamma = v/\allowbreak{}Lambda): the terms linear in (h, r) are total derivatives, and to second order L == (1/\allowbreak{}2)(dh)\textasciicircum{}2 + (1/\allowbreak{}2)(1 + 6 xi gamma\textasciicircum{}2)(dr)\textasciicircum{}2 + 6 xi gamma dh.dr - (1/\allowbreak{}2) m\_\allowbreak{}h0\textasciicircum{}2 h\textasciicircum{}2 - (1/\allowbreak{}2) m\_\allowbreak{}r0\textasciicircum{}2 r\textasciicircum{}2 up to total derivatives (Euler-Lagrange expressions in x0 and x1); a2 drops out} & True & 26.6 \\
31 & \texttt{canonical kinetic terms: Transpose[S].K.S == 1 for S = \{\{1, -6 xi gamma/\allowbreak{}Z\}, \{0, 1/\allowbreak{}Z\}\} if and only if Z\textasciicircum{}2 = 1 + 6 xi gamma\textasciicircum{}2 (1 - 6 xi) (the note's Z); with the note's literal kinetic terms (no 6 xi gamma\textasciicircum{}2 for the radion) it would be Z\textasciicircum{}2 = 1 - 36 xi\textasciicircum{}2 gamma\textasciicircum{}2} & True & 26.6 \\
32 & \texttt{mass mixing: with (h', r') = Q(theta).(h\_\allowbreak{}m, r\_\allowbreak{}m), Q = \{\{cos, -sin\}, \{sin, cos\}\}, the off-diagonal element of Transpose[Q].M'.Q is -(1/\allowbreak{}2)(M'\_\allowbreak{}11 - M'\_\allowbreak{}22) sin 2 theta + M'\_\allowbreak{}12 cos 2 theta, zero for tan 2 theta = 2 M'\_\allowbreak{}12/\allowbreak{}(M'\_\allowbreak{}11 - M'\_\allowbreak{}22) = 12 gamma xi Z m\_\allowbreak{}h0\textasciicircum{}2/\allowbreak{}(m\_\allowbreak{}r0\textasciicircum{}2 - m\_\allowbreak{}h0\textasciicircum{}2 (Z\textasciicircum{}2 - 36 xi\textasciicircum{}2 gamma\textasciicircum{}2)) (the note's formula); with the inverse matrix (the note's display) it is +(1/\allowbreak{}2)(M'\_\allowbreak{}11 - M'\_\allowbreak{}22) sin 2 theta + M'\_\allowbreak{}12 cos 2 theta (theta -> -theta); (h, r) = S.Q.(h\_\allowbreak{}m, r\_\allowbreak{}m) gives C\_\allowbreak{}h = cos theta - (6 xi gamma/\allowbreak{}Z) sin theta, C\_\allowbreak{}r = sin theta/\allowbreak{}Z} & True & 26.6 \\
33 & \texttt{Higgs-radion mixing in 7+1 (gamma = v/\allowbreak{}(Sqrt[3] M\_\allowbreak{}Pl), xi = 1, m\_\allowbreak{}r0 = 10 GeV, m\_\allowbreak{}h0 = 125.20 GeV; 60-digit arithmetic): |Z - 1| < 10\textasciicircum{}-30, |theta| = 3.5 10\textasciicircum{}-16, the rates of h -> gamma gamma and h -> g g change by less than 10\textasciicircum{}-30; in RS1 with <mu> = 3 TeV gamma = v/\allowbreak{}(3 TeV) = 0.082} & True & 26.6 \\*
\end{longtable}}
\Needspace{8\baselineskip}\subsection*{7.2 The independent Python program}\addcontentsline{toc}{subsection}{7.2 The independent Python program}
tools/ref\_radion.py — its boolean checks:

{\footnotesize
\begin{longtable}{>{\raggedright\arraybackslash}p{\dimexpr 0.0550\linewidth-2\tabcolsep\relax}>{\raggedright\arraybackslash}p{\dimexpr 0.8250\linewidth-2\tabcolsep\relax}>{\raggedright\arraybackslash}p{\dimexpr 0.1200\linewidth-2\tabcolsep\relax}}
\toprule
\textbf{\#} & \textbf{check} & \textbf{result} \\
\midrule\endfirsthead
\toprule
\textbf{\#} & \textbf{check} & \textbf{result} \\
\midrule\endhead
\endfoot
\bottomrule\endlastfoot
1 & \texttt{the 4-sphere of radius r0: Ricci\_\allowbreak{}ij = 3 gamma\_\allowbreak{}ij, off-diagonal 0, R = -12/\allowbreak{}r0\textasciicircum{}2 in the signature of 7+1.nb; volume of the unit S\textasciicircum{}4 = 8 pi\textasciicircum{}2/\allowbreak{}3} & True \\
2 & \texttt{radion on T\textasciicircum{}4: the reduced action equals -(M\_\allowbreak{}Pl\textasciicircum{}2/\allowbreak{}2) R\_\allowbreak{}E + 6 M\_\allowbreak{}Pl\textasciicircum{}2 (d beta)\textasciicircum{}2 - e\textasciicircum{}(-4 beta)(M\_\allowbreak{}Pl\textasciicircum{}2 Lambda + V0 U) up to a total derivative (Euler-Lagrange for N, a, beta), M\_\allowbreak{}Pl\textasciicircum{}2 = V0/\allowbreak{}kappa} & True \\
3 & \texttt{radion on S\textasciicircum{}4: the same with the potential -(M\_\allowbreak{}Pl\textasciicircum{}2/\allowbreak{}2) e\textasciicircum{}(-4 beta) 12/\allowbreak{}(r0\textasciicircum{}2 e\textasciicircum{}(2 beta)) from the curvature of the sphere, V0 = (8 pi\textasciicircum{}2/\allowbreak{}3) r0\textasciicircum{}4} & True \\
4 & \texttt{n extra dimensions (T\textasciicircum{}n, n = 1..4): kinetic term (M\_\allowbreak{}Pl\textasciicircum{}2/\allowbreak{}2)(n(n+2)/\allowbreak{}2)(dB)\textasciicircum{}2; the note's beta\_\allowbreak{}n = sqrt(2/\allowbreak{}(n(n+2))) gives (M\_\allowbreak{}Pl\textasciicircum{}2/\allowbreak{}2)(d phi)\textasciicircum{}2, not (M\_\allowbreak{}Pl\textasciicircum{}2/\allowbreak{}2)(1/\allowbreak{}2)(d phi)\textasciicircum{}2; beta\_\allowbreak{}n = 1/\allowbreak{}sqrt(n(n+2)) gives the 1/\allowbreak{}2} & True \\
5 & \texttt{the note's eight-dimensional alpha = -1/\allowbreak{}sqrt(3), beta = 1/\allowbreak{}sqrt(12): 2 alpha + 4 beta = 0, 4 beta = 2/\allowbreak{}sqrt(3), sqrt(12) beta = 1} & True \\
6 & \texttt{the kinetic coefficient solved from the reduction (n = 1..4) is n(n+2)/\allowbreak{}2} & True \\
7 & \texttt{the cosmological metric with the sphere of radius b: G\textasciicircum{}A\_\allowbreak{}A = torus + (6, 6, 6, 6, 3, 3, 3, 3)/\allowbreak{}b\textasciicircum{}2, R = torus - 12/\allowbreak{}b\textasciicircum{}2} & True \\
8 & \texttt{radion equation on S\textasciicircum{}4: Hb' + 3 Ha Hb + 4 Hb\textasciicircum{}2 = kappa (eps/\allowbreak{}6 - pa/\allowbreak{}2 + pb/\allowbreak{}3) + Lambda/\allowbreak{}3 - 3/\allowbreak{}b\textasciicircum{}2} & True \\
9 & \texttt{radion on T\textasciicircum{}4 with U: flat vacuum Lambda + kappa U0 = 0, U1 = 0; m\_\allowbreak{}r\textasciicircum{}2 = kappa K\_\allowbreak{}U/\allowbreak{}12} & True \\
10 & \texttt{radion on S\textasciicircum{}4 with U: Lambda + kappa U0 = 6/\allowbreak{}r0\textasciicircum{}2, kappa U1 = -12/\allowbreak{}r0\textasciicircum{}2, m\_\allowbreak{}r\textasciicircum{}2 = kappa K\_\allowbreak{}U/\allowbreak{}12 - 2/\allowbreak{}r0\textasciicircum{}2, from V and from the 8D radion equation} & True \\
11 & \texttt{S\textasciicircum{}4 with Lambda only: dS4 x S\textasciicircum{}4 (r0\textasciicircum{}2 = 9/\allowbreak{}Lambda, H\textasciicircum{}2 = 1/\allowbreak{}r0\textasciicircum{}2) solves the 00, 11, 44 equations; maximum of V with m\textasciicircum{}2 = -6 H\textasciicircum{}2, also from the 8D radion equation} & True \\
12 & \texttt{S\textasciicircum{}4 with a 4-form flux: eps = rho\_\allowbreak{}F, pa = -rho\_\allowbreak{}F, pb = +rho\_\allowbreak{}F; Minkowski needs Lambda = 9/\allowbreak{}(2 r0\textasciicircum{}2), kappa rho\_\allowbreak{}F = 3/\allowbreak{}(2 r0\textasciicircum{}2); m\textasciicircum{}2 = 6/\allowbreak{}r0\textasciicircum{}2 from V and from the 8D radion equation} & True \\
13 & \texttt{the note's Goldberger-Wise potential: for a static Phi(y) the Einstein-frame energy carries e\textasciicircum{}(-4 beta) (gradient part e\textasciicircum{}(-6 beta)), not e\textasciicircum{}(+4 beta)} & True \\
14 & \texttt{radion coupling to a mass on the brane: -M e\textasciicircum{}(-n beta/\allowbreak{}2) = -M + (n/\allowbreak{}2) M beta; coupling sqrt(n/\allowbreak{}(2(n+2)))/\allowbreak{}M\_\allowbreak{}Pl = 1/\allowbreak{}(sqrt(3) M\_\allowbreak{}Pl) for n = 4; fifth force 2 (coupling M\_\allowbreak{}Pl)\textasciicircum{}2 = n/\allowbreak{}(n+2) = 2/\allowbreak{}3} & True \\
15 & \texttt{Eot-Wash: m\_\allowbreak{}r > hbar c/\allowbreak{}38.6 micrometres = 5.1 meV; range of a 10 GeV radion 2.0e-17 m; gamma = v/\allowbreak{}(sqrt(3) M\_\allowbreak{}Pl) = 5.8e-17} & True \\
16 & \texttt{zonal harmonics on S\textasciicircum{}4 (l = 0..6): eigenvalue -l(l+3), C\_\allowbreak{}l(1) = (l+1)(l+2)/\allowbreak{}2, Int (C\_\allowbreak{}l/\allowbreak{}C\_\allowbreak{}l(1))\textasciicircum{}2 = V/\allowbreak{}d\_\allowbreak{}l; d\_\allowbreak{}l = binomial(l+4,4) - binomial(l+2,4) (l <= 20)} & True \\
17 & \texttt{Weyl's law on S\textasciicircum{}4: N(x)/\allowbreak{}(x\textasciicircum{}4/\allowbreak{}12) - 1 = -1/\allowbreak{}x\textasciicircum{}2 at x = 10, 100, 1000} & True \\
18 & \texttt{the brane at the north pole of S\textasciicircum{}4: the embedding induces r0\textasciicircum{}2 gamma; north pole (0,0,0,0,r0); 6 of the 10 rotations of R\textasciicircum{}5 fix it (SO(4)), 4 move it} & True \\
19 & \texttt{Killing spinors of S\textasciicircum{}4: nabla eps\_\allowbreak{}s = s (i/\allowbreak{}(2 r0)) gt eps\_\allowbreak{}s and D eps\_\allowbreak{}s = (2 s i/\allowbreak{}r0) eps\_\allowbreak{}s (s = +-1; residual < 1e-25 at four random points, 30 digits); Friedrich bound 4/\allowbreak{}r0\textasciicircum{}2} & True \\
20 & \texttt{Dirac spectrum on S\textasciicircum{}4: Weyl's law for spinors L\textasciicircum{}4/\allowbreak{}3 versus the degeneracies 2 x 4 binomial(l+3,3) (2.0e-3 at L = 1000)} & True \\
21 & \texttt{the sphere with the volume of the torus today: r = 688.9 GeV\textasciicircum{}-1; lowest levels 4.03 MeV (torus) and 2.90 MeV (sphere); dirac16complex on the sphere 2.90 MeV; levels below 5 GeV agree with Weyl} & True \\
22 & \texttt{RS1 (5D, mostly minus): G\textasciicircum{}mu\_\allowbreak{}nu = 3 A'' - 6 A'\textasciicircum{}2, G\textasciicircum{}y\_\allowbreak{}y = -6 A'\textasciicircum{}2; bulk Lambda = -6 k\textasciicircum{}2; tensions f\_\allowbreak{}UV = +6k/\allowbreak{}kappa, f\_\allowbreak{}IR = -6k/\allowbreak{}kappa} & True \\
23 & \texttt{Israel junction at the UV brane: K\textasciicircum{}mu\_\allowbreak{}nu = -A' delta, [K\textasciicircum{}mu\_\allowbreak{}nu] - delta [K] = 6k = 3 x (jump of A'), the delta part of G\textasciicircum{}mu\_\allowbreak{}nu; in the signature (-,+,+,+,+) [K - g K] = -kappa S} & True \\
24 & \texttt{RS1: M\_\allowbreak{}Pl\textasciicircum{}2 = (1/\allowbreak{}kappa) Int e\textasciicircum{}(-2k|y|) dy over [-pi R, pi R] = (1 - e\textasciicircum{}(-2 pi k R))/\allowbreak{}(kappa k)} & True \\
25 & \texttt{RS1 radion: kinetic term (3/\allowbreak{}(kappa k)) (d e\textasciicircum{}(-k pi T))\textasciicircum{}2, canonical mu = sqrt(6/\allowbreak{}(kappa k)) e\textasciicircum{}(-k pi T)} & True \\
26 & \texttt{8D warped along one extra direction: G\textasciicircum{}y\_\allowbreak{}y = -6 A'\textasciicircum{}2, G\textasciicircum{}z\_\allowbreak{}z = 4 A'' - 10 A'\textasciicircum{}2; with only Lambda A'' = 0 and A' = 0} & True \\
27 & \texttt{Goldberger-Wise: Phi'' - 4k Phi' - m\textasciicircum{}2 Phi = 0 from the bulk action, solved by e\textasciicircum{}(2ky)(C1 e\textasciicircum{}(nu k y) + C2 e\textasciicircum{}(-nu k y)), nu = sqrt(4 + m\textasciicircum{}2/\allowbreak{}k\textasciicircum{}2) = 2 + eps; V'(<mu>) = 0 at (<mu>/\allowbreak{}f)\textasciicircum{}eps = v\_\allowbreak{}B/\allowbreak{}v\_\allowbreak{}A, V''(<mu>) = 8 k eps\textasciicircum{}2 v\_\allowbreak{}B\textasciicircum{}2 <mu>\textasciicircum{}2/\allowbreak{}f\textasciicircum{}4; m\_\allowbreak{}r\textasciicircum{}2 = (4/\allowbreak{}3) kappa k\textasciicircum{}2 eps\textasciicircum{}2 v\_\allowbreak{}B\textasciicircum{}2 (<mu>/\allowbreak{}f)\textasciicircum{}2 (note: 2/\allowbreak{}3)} & True \\
28 & \texttt{conformally flat h = Omega\textasciicircum{}2 eta: sqrt|h| R(h) = -6 Omega box Omega, box = d0\textasciicircum{}2 - d1\textasciicircum{}2} & True \\
29 & \texttt{the Higgs boson on the brane: linear terms are total derivatives; to second order L = (1/\allowbreak{}2)(dh)\textasciicircum{}2 + (1/\allowbreak{}2)(1 + 6 xi gamma\textasciicircum{}2)(dr)\textasciicircum{}2 + 6 xi gamma dh.dr - (1/\allowbreak{}2) m\_\allowbreak{}h0\textasciicircum{}2 h\textasciicircum{}2 - (1/\allowbreak{}2) m\_\allowbreak{}r0\textasciicircum{}2 r\textasciicircum{}2 up to total derivatives; o\_\allowbreak{}2 drops out} & True \\
30 & \texttt{the shift h = h' - (6 xi gamma/\allowbreak{}Z) r', r = r'/\allowbreak{}Z is canonical iff Z\textasciicircum{}2 = 1 + 6 xi gamma\textasciicircum{}2 (1 - 6 xi) (note's literal kinetic terms: 1 - 36 xi\textasciicircum{}2 gamma\textasciicircum{}2); tan 2 theta = 12 gamma xi Z m\_\allowbreak{}h0\textasciicircum{}2/\allowbreak{}(m\_\allowbreak{}r0\textasciicircum{}2 - m\_\allowbreak{}h0\textasciicircum{}2 (Z\textasciicircum{}2 - 36 xi\textasciicircum{}2 gamma\textasciicircum{}2)) diagonalizes with (h', r') = Q(theta)(h\_\allowbreak{}m, r\_\allowbreak{}m); the inverse matrix needs -theta; C\_\allowbreak{}h = cos - (6 xi gamma/\allowbreak{}Z) sin, C\_\allowbreak{}r = sin/\allowbreak{}Z} & True \\
31 & \texttt{Higgs-radion mixing in 7+1 (gamma = v/\allowbreak{}(sqrt(3) M\_\allowbreak{}Pl), xi = 1, m\_\allowbreak{}r0 = 10 GeV): |Z - 1| < 1e-30, |theta| = 3.5e-16, rates change by < 1e-30; warped gamma = v/\allowbreak{}(3 TeV) = 0.082} & True \\*
\end{longtable}}
\textbf{Numbers printed by both programs} (largest relative difference per group, and the tolerance used by tools/verify\_7p1\_nb.py; integers must agree exactly):

{\footnotesize
\begin{longtable}{>{\raggedright\arraybackslash}p{\dimexpr 0.5200\linewidth-2\tabcolsep\relax}>{\raggedright\arraybackslash}p{\dimexpr 0.1000\linewidth-2\tabcolsep\relax}>{\raggedright\arraybackslash}p{\dimexpr 0.2200\linewidth-2\tabcolsep\relax}>{\raggedright\arraybackslash}p{\dimexpr 0.1600\linewidth-2\tabcolsep\relax}}
\toprule
\textbf{quantity} & \textbf{values} & \textbf{max. relative difference} & \textbf{tolerance} \\
\midrule\endfirsthead
\toprule
\textbf{quantity} & \textbf{values} & \textbf{max. relative difference} & \textbf{tolerance} \\
\midrule\endhead
\endfoot
\bottomrule\endlastfoot
kinetic coefficient $n(n+2)/2$, $n=1,\dots,4$ (solved from the reduction) & 4 & 0 & exact \\
$\beta_n^2$ of the note and the canonical $\beta_n^2$, $n=4$ & 2 & 0 & $1\times10^{-15}$ \\
$S^4$ with $U$: $(\Lambda+\kappa U_0)r_0^2$, $\kappa U_1r_0^2$, $(m_r^2-\kappa K_U/12)r_0^2$ & 3 & 0 & exact \\
$S^4$ with $\Lambda$ only: $\Lambda r_0^2$, $H^2r_0^2$, $m^2/H^2$ & 3 & 0 & exact \\
$S^4$ with a 4-form flux: $\Lambda r_0^2$, $\kappa\rho_Fr_0^2$, $m^2r_0^2$ & 3 & 0 & $1\times10^{-15}$ \\
E\textbackslash{}"ot-Wash bound, range of a 10 GeV radion, $\gamma=v/(\sqrt3M_{\rm Pl})$ & 3 & $2.1\times10^{-16}$ & $1\times10^{-10}$ \\
multiplicities $d_l$ on $S^4$, $l=0,\dots,8$ & 9 & 0 & exact \\
Weyl's law on $S^4$: $N/(x^4/12)-1$ at $x=10,100,1000$ & 3 & $8.2\times10^{-11}$ & $1\times10^{-9}$ \\
Dirac degeneracies on $S^4$ and the spinor Weyl ratio & 7 & $7.0\times10^{-14}$ & $1\times10^{-9}$ \\
today: $r_{S^4}$, lowest masses, numbers of levels below 5 GeV & 8 & $2.2\times10^{-16}$ & $1\times10^{-12}$ \\
RS1: $3$, $f_{\rm UV}\kappa/k$, $f_{\rm IR}\kappa/k$, $f^2\kappa k$ & 4 & 0 & exact \\
strength of the warped coupling over the flat one & 1 & $1.8\times10^{-16}$ & $1\times10^{-12}$ \\
Goldberger--Wise mass coefficient: $\frac43$ and the note's $\frac23$ & 2 & 0 & $1\times10^{-15}$ \\
Higgs--radion: $c=-6$, $K_{hr}/(\xi\gamma)$, $(K_{rr}-1)/(\xi\gamma^2)$ & 3 & 0 & exact \\
Higgs--radion mixing in 7+1: $Z-1$, $\theta$, rate changes; warped $\gamma$ & 5 & $4.4\times10^{-16}$ & $1\times10^{-12}$ \\*
\end{longtable}}
\Needspace{8\baselineskip}\subsection*{7.3 The earlier sections are unchanged}\addcontentsline{toc}{subsection}{7.3 The earlier sections are unchanged}
tools/verify\_7p1\_nb.py compares the evaluated notebook with the backup of section 4 and with the Python program; its lines on section 26:

\begin{Verbatim}[fontsize=\small,frame=single,framesep=2mm]
section 26: check outputs (all True) : 33 (True)
stage 7: without section 26 (and with the old title of the summary section) the 815 input and
    ...text cells equal those of the backup taken before section 26 was appended : True
stage 7: the 515 outputs of sections 0-25 and of the summary equal those of the backup (summary
    ...counts and tables, the private context name and measured times excepted) : True
section 26 numbers: 15 groups of printed numbers agree with tools/ref_radion.py (integers and
    ...exact values exactly, floating-point numbers to 1e-15 ... 1e-9 per group) : True
VERDICT: ALL CROSS-CHECKS PASSED
\end{Verbatim}
\Needspace{8\baselineskip}\subsection*{7.4 Files}\addcontentsline{toc}{subsection}{7.4 Files}
\textbf{SHA-256 of the files} (at the time this document was generated):

\begin{Verbatim}[fontsize=\small,frame=single,framesep=2mm]
7+1.nb                             41d252f9852f210c228f45e1ac956d2d7ef26faeb0062a5880280e1aae0ef
    ...fc7
tools/make_7p1_nb.py               996d26c9552cd44902560af16a1683fa71457261f18fefaeddae87bd99a2b
    ...144
tools/nb_part_gravity.py           c4caf2ad5c932e8f603341edade5f4b98fa784feb05906c051eab7c884cd6
    ...8f8
tools/nb_part_ks_gravity.py        a656c0ad040b4576b75b17e029545109e531c8e656ad99bfdb5cba946e967
    ...772
tools/nb_part_sm.py                43e92beed57453a4b57931de6d603ac6c218efba6e4005133b69b8054a975
    ...c5b
tools/nb_part_coupled.py           e630521d64132e71afc861114f699eb06454557cf14a57ad8ec1eef96247a
    ...ae0
tools/nb_part_dark.py              10ad8fb8ae25fa07bbee758c653995acd4e2c80733059a8b3648bd1593e46
    ...db0
tools/nb_part_torus_brane.py       fb904ff8d85b904c55b3b2a8f60fb0cb047f1ccf1b98db3364092f2bae0b2
    ...162
tools/nb_part_radion.py            53c5604707a3e22228c840ee0a2d2125616f9410f75a5535041dfaf18f497
    ...c00
tools/lovelock71.py                d93ecbd0f82adb972218d59dfe32999a4af944945ca9758607f3a633e1af9
    ...5a1
tools/ref_gravity.py               f07cf8adaff585af89c87d19532f1aaa7db606ad88ae90483b2d1f41341ed
    ...f36
tools/ref_ks_gravity.py            b6488ef24b1c16aacd8ff6cf6a5a1ba3f316e082449ad97480022d1a26383
    ...27f
tools/ref_sm.py                    f3f80806893f42bbf166e1f1de896ffd611dbfe847685da8e3ee367e9cde6
    ...aba
tools/ref_coupled.py               695dbd9375753653472d9c37502f0067c4a64fe01a64ba8c05e8b2c5da85a
    ...a58
tools/ref_dark.py                  1600e6b57ab0383b4f2b664e2796d3864c9afa77a9ed59f007f66e1139453
    ...ff5
tools/ref_torus_brane.py           1f623656ab0f5c8ee43cb310f7bfd580cbba5a7385da2963adaae4d2f8bda
    ...af7
tools/ref_radion.py                e9d47176b9f2d2232a84b3b82188bd3b0a19cde3dbe1ba013c1fe5abff634
    ...1ce
tools/ref71.py                     1a7a06eb0add3f459611776dc9c474a5bb959d947a9db949db4612b1ab943
    ...4e0
tools/ref44.py                     ab100c3ebcc055fbe59c132a02f80ecb17630e934e99d59f88e919ec8064f
    ...af5
\end{Verbatim}
\textbf{Output of tools/ref\_radion.py (independent checks of section 26; exit code 0)} (complete):

\begin{Verbatim}[fontsize=\small,frame=single,framesep=2mm]
the 4-sphere of radius r0: Ricci_ij = 3 gamma_ij, off-diagonal 0, R = -12/r0^2 in the signature
    ...of 7+1.nb; volume of the unit S^4 = 8 pi^2/3 : True
radion on T^4: the reduced action equals -(M_Pl^2/2) R_E + 6 M_Pl^2 (d beta)^2 - e^(-4
    ...beta)(M_Pl^2 Lambda + V0 U) up to a total derivative (Euler-Lagrange for N, a, beta),
    ...M_Pl^2 = V0/kappa : True
radion on S^4: the same with the potential -(M_Pl^2/2) e^(-4 beta) 12/(r0^2 e^(2 beta)) from
    ...the curvature of the sphere, V0 = (8 pi^2/3) r0^4 : True
n extra dimensions (T^n, n = 1..4): kinetic term (M_Pl^2/2)(n(n+2)/2)(dB)^2; the note's beta_n
    ...= sqrt(2/(n(n+2))) gives (M_Pl^2/2)(d phi)^2, not (M_Pl^2/2)(1/2)(d phi)^2; beta_n =
    ...1/sqrt(n(n+2)) gives the 1/2 : True
the note's eight-dimensional alpha = -1/sqrt(3), beta = 1/sqrt(12): 2 alpha + 4 beta = 0, 4
    ...beta = 2/sqrt(3), sqrt(12) beta = 1 : True
the kinetic coefficient solved from the reduction (n = 1..4) is n(n+2)/2 : True
NUM general n: kinetic coefficient n(n+2)/2 for n = 1..4 : (1.5, 4.0, 7.5, 12.0)
NUM the note's beta_n^2 = 2/(n(n+2)) and the canonical 1/(n(n+2)), n = 4 :
    ...(0.08333333333333333, 0.041666666666666664)
the cosmological metric with the sphere of radius b: G^A_A = torus + (6, 6, 6, 6, 3, 3, 3,
    ...3)/b^2, R = torus - 12/b^2 : True
radion equation on S^4: Hb' + 3 Ha Hb + 4 Hb^2 = kappa (eps/6 - pa/2 + pb/3) + Lambda/3 - 3/b^2
    ...: True
radion on T^4 with U: flat vacuum Lambda + kappa U0 = 0, U1 = 0; m_r^2 = kappa K_U/12 : True
radion on S^4 with U: Lambda + kappa U0 = 6/r0^2, kappa U1 = -12/r0^2, m_r^2 = kappa K_U/12 -
    ...2/r0^2, from V and from the 8D radion equation : True
S^4 with Lambda only: dS4 x S^4 (r0^2 = 9/Lambda, H^2 = 1/r0^2) solves the 00, 11, 44
    ...equations; maximum of V with m^2 = -6 H^2, also from the 8D radion equation : True
S^4 with a 4-form flux: eps = rho_F, pa = -rho_F, pb = +rho_F; Minkowski needs Lambda = 9/(2
    ...r0^2), kappa rho_F = 3/(2 r0^2); m^2 = 6/r0^2 from V and from the 8D radion equation :
    ...True
the note's Goldberger-Wise potential: for a static Phi(y) the Einstein-frame energy carries
    ...e^(-4 beta) (gradient part e^(-6 beta)), not e^(+4 beta) : True
NUM sphere with U: Lambda + kappa U0, kappa U1, m_r^2 - kappa K_U/12 (units 1/r0^2) : (6.0,
    ...-12.0, -2.0)
NUM sphere with Lambda only: Lambda r0^2, H^2 r0^2, m^2/H^2 : (9.0, 1.0, -6.0)
NUM sphere with a 4-form flux: Lambda r0^2, kappa rho_F r0^2, m^2 r0^2 : (4.5, 1.5, 6.0)
radion coupling to a mass on the brane: -M e^(-n beta/2) = -M + (n/2) M beta; coupling
    ...sqrt(n/(2(n+2)))/M_Pl = 1/(sqrt(3) M_Pl) for n = 4; fifth force 2 (coupling M_Pl)^2 =
    ...n/(n+2) = 2/3 : True
Eot-Wash: m_r > hbar c/38.6 micrometres = 5.1 meV; range of a 10 GeV radion 2.0e-17 m; gamma =
    ...v/(sqrt(3) M_Pl) = 5.8e-17 : True
NUM Eot-Wash bound on the radion mass (eV), range of a 10 GeV radion (m), gamma = v/(sqrt(3)
    ...M_Pl) : (0.005112097937823834, 1.973269804e-17, 5.837986676464508e-17)
zonal harmonics on S^4 (l = 0..6): eigenvalue -l(l+3), C_l(1) = (l+1)(l+2)/2, Int
    ...(C_l/C_l(1))^2 = V/d_l; d_l = binomial(l+4,4) - binomial(l+2,4) (l <= 20) : True
Weyl's law on S^4: N(x)/(x^4/12) - 1 = -1/x^2 at x = 10, 100, 1000 : True
NUM S^4 multiplicities d_l, l = 0..8 : (1, 5, 14, 30, 55, 91, 140, 204, 285)
NUM S^4 Weyl: N(x)/(x^4/12) - 1 at x = mu r = 10, 100, 1000 : (-0.010000000000000009,
    ...-9.999999999998899e-05, -9.999999999177334e-07)
the brane at the north pole of S^4: the embedding induces r0^2 gamma; north pole (0,0,0,0,r0);
    ...6 of the 10 rotations of R^5 fix it (SO(4)), 4 move it : True
Killing spinors of S^4: nabla eps_s = s (i/(2 r0)) gt eps_s and D eps_s = (2 s i/r0) eps_s (s =
    ...+-1; residual < 1e-25 at four random points, 30 digits); Friedrich bound 4/r0^2 : True
Dirac spectrum on S^4: Weyl's law for spinors L^4/3 versus the degeneracies 2 x 4
    ...binomial(l+3,3) (2.0e-3 at L = 1000) : True
NUM Dirac on S^4: degeneracy 4 binomial(l + 3, 3) of +i(l + 2)/r, l = 0..5; Weyl ratio - 1 at L
    ...= 1000 : (4, 16, 40, 80, 140, 224, 0.0019989980000001406)
the sphere with the volume of the torus today: r = 688.9 GeV^-1; lowest levels 4.03 MeV (torus)
    ...and 2.90 MeV (sphere); dirac16complex on the sphere 2.90 MeV; levels below 5 GeV agree
    ...with Weyl : True
NUM today: r_S4 (GeV^-1, m), lowest masses (MeV): torus 2 pi/l_perp, sphere scalar 2/r, sphere
    ...dirac16complex 2/r; levels below 5 GeV: torus lattice, sphere exact, Weyl :
    ...(688.9415270913784, 1.359467512131065e-13, 4.026524867292262, 2.903003987063663,
    ...2.903003987063663, 11733511217161, 11737500461050, 11733517482065.238)
RS1 (5D, mostly minus): G^mu_nu = 3 A'' - 6 A'^2, G^y_y = -6 A'^2; bulk Lambda = -6 k^2;
    ...tensions f_UV = +6k/kappa, f_IR = -6k/kappa : True
Israel junction at the UV brane: K^mu_nu = -A' delta, [K^mu_nu] - delta [K] = 6k = 3 x (jump of
    ...A'), the delta part of G^mu_nu; in the signature (-,+,+,+,+) [K - g K] = -kappa S : True
RS1: M_Pl^2 = (1/kappa) Int e^(-2k|y|) dy over [-pi R, pi R] = (1 - e^(-2 pi k R))/(kappa k) :
    ...True
RS1 radion: kinetic term (3/(kappa k)) (d e^(-k pi T))^2, canonical mu = sqrt(6/(kappa k))
    ...e^(-k pi T) : True
8D warped along one extra direction: G^y_y = -6 A'^2, G^z_z = 4 A'' - 10 A'^2; with only Lambda
    ...A'' = 0 and A' = 0 : True
NUM RS1: coefficient of A'' in G^0_0, f_UV kappa/k, f_IR kappa/k, f^2 kappa k (canonical
    ...radion) : (3.0, 6.0, -6.0, 6.0)
NUM strength of the radion coupling: warped with <mu> = 3 TeV over 7+1, i.e. sqrt(3) M_Pl/(3
    ...TeV) : 1405847905476738.8
Goldberger-Wise: Phi'' - 4k Phi' - m^2 Phi = 0 from the bulk action, solved by e^(2ky)(C1 e^(nu
    ...k y) + C2 e^(-nu k y)), nu = sqrt(4 + m^2/k^2) = 2 + eps; V'(<mu>) = 0 at (<mu>/f)^eps =
    ...v_B/v_A, V''(<mu>) = 8 k eps^2 v_B^2 <mu>^2/f^4; m_r^2 = (4/3) kappa k^2 eps^2 v_B^2
    ...(<mu>/f)^2 (note: 2/3) : True
NUM Goldberger-Wise: m_r^2/(kappa k^2 eps^2 v_B^2 (<mu>/f)^2) with f^2 = 6/(kappa k) and with
    ...the note's f^2 = 12/(kappa k) : (1.3333333333333333, 0.6666666666666666)
conformally flat h = Omega^2 eta: sqrt|h| R(h) = -6 Omega box Omega, box = d0^2 - d1^2 : True
the Higgs boson on the brane: linear terms are total derivatives; to second order L =
    ...(1/2)(dh)^2 + (1/2)(1 + 6 xi gamma^2)(dr)^2 + 6 xi gamma dh.dr - (1/2) m_h0^2 h^2 -
    ...(1/2) m_r0^2 r^2 up to total derivatives; o_2 drops out : True
the shift h = h' - (6 xi gamma/Z) r', r = r'/Z is canonical iff Z^2 = 1 + 6 xi gamma^2 (1 - 6
    ...xi) (note's literal kinetic terms: 1 - 36 xi^2 gamma^2); tan 2 theta = 12 gamma xi Z
    ...m_h0^2/(m_r0^2 - m_h0^2 (Z^2 - 36 xi^2 gamma^2)) diagonalizes with (h', r') =
    ...Q(theta)(h_m, r_m); the inverse matrix needs -theta; C_h = cos - (6 xi gamma/Z) sin, C_r
    ...= sin/Z : True
Higgs-radion mixing in 7+1 (gamma = v/(sqrt(3) M_Pl), xi = 1, m_r0 = 10 GeV): |Z - 1| < 1e-30,
    ...|theta| = 3.5e-16, rates change by < 1e-30; warped gamma = v/(3 TeV) = 0.082 : True
NUM Higgs-radion mixing: c in sqrt|h| R(h) = c Omega box Omega, K_hr/(xi gamma), (K_rr - 1)/(xi
    ...gamma^2) : (-6.0, 6.0, 6.0)
NUM Higgs-radion mixing in 7+1 (gamma = v/(sqrt(3) M_Pl), xi = 1, m_r0 = 10 GeV): Z - 1, theta,
    ...Gamma(h -> gamma gamma)/SM - 1, Gamma(h -> g g)/SM - 1; warped gamma = v/(3 TeV) :
    ...(-5.112313265186566e-32, -3.525281784433872e-16, 9.947138056858168e-32,
    ...3.387862153903166e-31, 0.08207321341308735)
boolean checks: 31   failed: 0
VERDICT: ALL CHECKS PASSED
\end{Verbatim}
\end{document}
